% Generated draft: requires mathematical and source review.
\documentclass[11pt]{article}
\usepackage[margin=25mm]{geometry}
\usepackage{fontspec}
\setmainfont{Latin Modern Roman}
\usepackage{amsmath,amssymb,mathtools,unicode-math}
\setmathfont{Latin Modern Math}
\usepackage{enumitem,needspace,xurl,hyperref}
\setlength{\parindent}{0pt}
\setlength{\parskip}{5pt}
\setlength{\emergencystretch}{3em}
\newcommand{\N}{\mathbb{N}}
\newcommand{\Z}{\mathbb{Z}}
\newcommand{\Q}{\mathbb{Q}}
\newcommand{\R}{\mathbb{R}}
\newcommand{\C}{\mathbb{C}}
\newcommand{\F}{\mathbb{F}}
\newcommand{\A}{\mathbb{A}}
\newcommand{\PP}{\mathbb{P}}
\newcommand{\E}{\mathbb{E}}
\newcommand{\eps}{\varepsilon}
\newcommand{\dd}{\,\mathrm{d}}
\begin{document}
\section*{UnsolvedMath ID: 701464}
\textbf{Source problem code:} catalogue entry
\subsection{Definitions and mathematical statement}
For a group \(H\), write \(H'=[H,H]\) and \(H''=[H',H']\). The group \(H\) is metabelian if \(H''=\{1\}\), equivalently if \(H'\) is abelian. Let \(F_2\) be the free group on two generators and let
\[
G=F_2/F_2'',
\]
the relatively free group of rank \(2\) in the variety of all metabelian groups.

For groups \(A,B\), a family \(\mathscr F\subseteq\operatorname{Hom}(A,B)\) is separating if
\[
\forall a\in A\setminus\{1\}\;\exists f\in\mathscr F:\quad f(a)\ne1.
\]
It is discriminating if
\[
\forall n\ge1\;\forall(a_1,\ldots,a_n)\in(A\setminus\{1\})^n\;\exists f\in\mathscr F:\quad
\bigwedge_{i=1}^{n}f(a_i)\ne1.
\]
For \(t\in T\), define its centralizer by
\[
C_T(t)=\{x\in T:xt=tx\}.
\]

\textbf{Question.} Let \(T\) be any metabelian group such that \(C_T(t)\) is abelian for every \(t\in T\setminus\{1\}\). Suppose that \(T\) admits a separating family of surjective homomorphisms \(T\to G\). Must \(T\) also admit a discriminating family of surjective homomorphisms \(T\to G\)? Equivalently, under these hypotheses, must every finite list \(t_1,\ldots,t_n\) of nonidentity elements of \(T\), with \(n\ge1\), admit a single surjective homomorphism \(\phi:T\to G\) such that \(\phi(t_i)\ne1\) for every \(i\)? No finite generation assumption on \(T\) is supplied.

The source states that a positive answer would give a metabelian analogue of a classical theorem of B. Baumslag; it does not specify that theorem or provide its citation.
\subsection{Short English statement}
Suppose a metabelian group \(T\) has abelian centralizers of all nonidentity elements, and each nonidentity element survives under some surjection onto the free metabelian group of rank \(2\). Must every finite collection of nonidentity elements survive simultaneously under one such surjection?
\subsection{Sources}
\begin{itemize}
\item Local UnsolvedMath registry, permanent ID \(701464\), source problem code \texttt{catalogue entry}, titled \emph{Kourovka Notebook Problem 18.35}. The supplied record labels the problem open; this label is source metadata, not an independently verified mathematical or literature assessment.
\item \emph{The Kourovka Notebook (Unsolved Problems in Group Theory)}, original problem 18.35, 18th issue (2014); the supplied background also identifies the 21st issue edition (2026), \texttt{arXiv:1401.0300v48}. Proposed by A. Gaglione and D. Spellman in 2014. The supplied collection description names E. I. Khukhro and V. D. Mazurov as editors.
\item The supplied statement attributes a related classical theorem to B. Baumslag, without a theorem statement or bibliographic reference.
\end{itemize}
\end{document}
